\chapter*{Lời nói đầu}
\addcontentsline{toc}{chapter}{Lời nói đầu}
Trong những năm gần đây, trí tuệ nhân tạo tạo sinh đã trở thành một trong các hướng công nghệ có tốc độ phát triển nhanh nhất. Các mô hình sinh ảnh, sinh văn bản và sinh video giúp người dùng tạo ra nội dung chất lượng cao trong thời gian rất ngắn. Tuy nhiên, khả năng tạo nội dung quá dễ dàng cũng làm tăng nguy cơ giả mạo thông tin, phát tán tin sai lệch, tạo bằng chứng hình ảnh không đáng tin cậy và làm suy giảm niềm tin vào dữ liệu số.

Báo cáo Project I này trình bày quá trình tìm hiểu và thiết kế nền tảng SourceVerify -- một hệ thống thử nghiệm nhằm hỗ trợ nhận diện dấu hiệu nội dung do AI tạo sinh. Nội dung báo cáo tập trung vào phương pháp tiếp cận đa tín hiệu, trong đó mỗi method phân tích một khía cạnh khác nhau của dữ liệu ảnh. Do phạm vi Project I chủ yếu là làm quen nền tảng công nghệ và cách đặt vấn đề, báo cáo không đặt mục tiêu chứng minh hệ thống đạt độ chính xác tuyệt đối, mà hướng tới việc trình bày rõ bài toán, cơ sở lý thuyết, thiết kế hệ thống và kết quả triển khai ban đầu.
